% ============================================================
% 第一章 引言
% ============================================================
\section{引言}

\subsection{研究背景与意义}

椭圆型偏微分方程是数学物理中最重要的方程类型之一，
广泛应用于流体力学、电磁场计算、声学传播、热传导等领域。
本文在统一框架下研究三类椭圆型方程：
\begin{equation}
  (-\nabla^2 + \sigma)u = f, \quad (x,y) \in \Omega
  \label{eq:unified_pde_intro}
\end{equation}
其中参数 $\sigma$ 的取值决定方程类型：
\begin{itemize}[itemsep=2pt]
  \item $\sigma = 0$：Poisson 方程 $-\nabla^2 u = f$，描述稳态热传导、静电势等物理过程；
  \item $\sigma = +\kappa^2$（$\kappa > 0$）：modified Helmholtz 方程（screened Poisson 方程），
        矩阵保持对称正定，条件数随 $\kappa^2$ 增大而改善；
  \item $\sigma = -\kappa^2$：true Helmholtz 方程（时谐波动方程的频域形式），
        矩阵不定，当特征值 $\lambda_{p,q}=\kappa^2$ 时出现共振。
\end{itemize}
Poisson 方程与 modified Helmholtz 方程的离散矩阵均为对称正定，
适用 CG 等基于正定性的迭代法；true Helmholtz 方程的离散矩阵不定，
需要 GMRES 或 MINRES 等更一般的 Krylov 方法~\cite{ShaoWenting2017, LiuJiayuan2019, LiSiqing2020}。

数值求解椭圆型方程的方法大致分为两类：
\textbf{直接法}和\textbf{迭代法}~\cite{MinTao2017}。
直接法通过矩阵分解或变换技术一次性得到精确解，
其典型代表包括基于快速 Fourier 变换（FFT）的谱方法；
迭代法则从初始猜测出发逐步逼近解，
其代表包括 Krylov 子空间方法如 GMRES。

在规则区域（如矩形域）上，FFT 直接法具有显著的效率优势：
利用离散正弦变换（DST）或离散余弦变换（DCT）的对角化性质，
可以将 $N^2$ 个耦合的差分方程转化为 $N^2$ 个独立的代数方程，
总计算量仅为 $O(N^2 \log N)$，远优于一般稀疏矩阵直接分解的 $O(N^3)$。
循环约化（CR）和 FACR 混合算法可进一步降低常数因子，
实际实现的复杂度为 $O(N^2 \log N)$。

然而，FFT 直接法的适用范围受限于区域形状和方程系数的可分离性。
对于非规则区域或变系数问题，迭代法（如 GMRES）更具灵活性，
但收敛速度受矩阵条件数影响，通常需要预处理技术加速。

因此，系统对比 FFT 直接法与 GMRES 迭代法在不同条件下的求解性能，
对于选择合适的数值方法具有重要的理论和实践意义。

\subsection{国内外研究现状}

\subsubsection{基于 FFT 的快速直接法}

Hockney \cite{Hockney1965} 首先提出了利用 Fourier 分析求解
离散 Poisson 方程的方法，计算复杂度为 $O(N^2 \log N)$。
Buzbee 等 \cite{Buzbee1970} 系统分析了循环约化（CR）方法，
并提出了 FACR 混合算法，在最优参数下达到 $O(N^2 \log\log N)$ 的复杂度。
Swarztrauber \cite{Swarztrauber1977} 对 CR、FA 和 FACR 方法
进行了完整的理论分析，明确了各种方法的理论复杂度。
Houstis 和 Papatheodorou \cite{HoustisPapatheodorou1979a, HoustisPapatheodorou1979b}
提出了基于九点紧致差分格式的 FFT9 方法，实现了四阶精度的快速求解。
在国内，邵文婷\cite{ShaoWenting2017}研究了不规则区域上 Poisson 方程的快速求解器，
李国燕等\cite{LiGuoyan2019}实现了基于 FFT 的泊松方程快速求解器的硬件加速。

在差分格式方面，Gupta 等 \cite{Gupta1984} 和 Spotz 等 \cite{Spotz1998}
发展了高阶紧致差分格式，Strang \cite{Strang1999} 对 DCT 的数学性质
做了精辟阐述，为 Neumann 边界条件下的 FFT 求解提供了理论基础。
朱爱玲\cite{ZhuAiling2015}研究了泊松方程五点差分格式的外推法，
进一步提高了低阶格式的精度。

\subsubsection{Helmholtz 方程求解}

Helmholtz 方程的数值求解需区分两种类型。
对于 modified Helmholtz 方程 $(-\nabla^2+\kappa^2)u=f$，
离散矩阵对称正定，频域分母 $\lambda_{p,q}+\kappa^2 > 0$ 恒成立，
条件数随 $\kappa^2$ 增大而改善，迭代法通常加速收敛。
对于 true Helmholtz 方程 $(-\nabla^2-\kappa^2)u=f$，
离散矩阵不定，当某特征值 $\lambda_{p,q}\approx\kappa^2$ 时频域分母接近零，
导致共振现象和数值困难~\cite{Bayliss1983, Ernst2011}。
Ernst 和 Gander~\cite{Ernst2011}系统分析了经典迭代法求解 true Helmholtz
方程的收敛性障碍，指出 shifted-Laplacian 预处理是缓解困难的主要手段。
Ihlenburg~\cite{Ihlenburg1998}对有限元方法求解 true Helmholtz 方程
进行了系统的误差分析。
国内方面，刘镓源\cite{LiuJiayuan2019}研究了两类 Helmholtz 问题的差分法，
姚登明\cite{YaoDengming2018}对带 PML 边界条件的 Helmholtz 方程
进行了有限差分方法研究。

\subsubsection{GMRES 迭代法}

Saad 和 Schultz \cite{SaadSchultz1986} 于 1986 年提出 GMRES 算法，
这是求解非对称线性系统的最重要的迭代法之一。
算法的核心思想是在 Krylov 子空间中极小化残差范数，
通过 Arnoldi 过程构建正交基，并用 Givens 旋转增量求解最小二乘问题。
Saad \cite{Saad2003} 对 Krylov 子空间方法做了系统总结。
牛强等\cite{Niu2009}提出了改进的 GMRES 方法，
有效提升了求解非对称线性方程组的收敛速度。

\subsection{本文工作与结构安排}

本文的主要工作和创新点如下：

\begin{enumerate}[itemsep=2pt]
  \item 在统一的 $(-\nabla^2+\sigma)u=f$ 框架下，系统实现了基于 FFT 的五种快速直接求解器
        （FA、CR、FACR、FFT9、5点基准），支持 Dirichlet、Neumann 和混合边界条件，
        并明确区分 modified Helmholtz（$\sigma=+\kappa^2$，正定）与 true Helmholtz（$\sigma=-\kappa^2$，不定）；
  \item 对 Neumann 边界条件，采用 ghost-point 方法配合对称化 DCT-I 变换，解决了非对称 Laplacian 矩阵的对角化问题；
  \item 实现了基于 Givens 旋转的 GMRES 迭代法，构造了 2D Helmholtz 稀疏矩阵，
        与 FFT 直接法进行了系统对比，并分析了 modified/true Helmholtz 方程下迭代行为的本质差异；
  \item 通过 Fourier 特征值分析证明了三参数九点修正族的六阶不可行性，
        为四阶紧致格式是最优选择提供了理论依据，并利用 Lean 4 进行了辅助形式化校验。
\end{enumerate}

本文结构安排如下：
第\ref{sec:math}章介绍数学基础与差分格式；
第\ref{sec:fft}章阐述基于 FFT 的快速直接法；
第\ref{sec:helmholtz}章讨论 Helmholtz 方程的求解与边界条件处理；
第\ref{sec:gmres}章介绍 GMRES 迭代法；
第\ref{sec:experiments}章给出数值实验结果；
第\ref{sec:conclusion}章总结全文并展望未来工作。
